\subsection{Ray-Tracing based Simulation}
Ray-tracing is a deterministic channel modeling technique that provides a physically grounded representation of radio signal propagation.
It models how electromagneticwaves interact with the enviroment by considering refelction, diffrations, and scattering.

The fundamental principle of ray-tracing is to approximate electromagnetic waves as rays whose paths are geometrically computed.
The channel repose between a transmitter and a receiveris derived by tracing multiple propagation of emitted rays which interact with surronding objects, undergoing reflections, diffractions, or scattering based on the material and geometry of enviroment.
Ray-tracing has been widely used sicne 1990s for modeling radio propagation in complex environments.
Its deterministic nature enable high accuracy in predicting signal strength, delay spread, and angle of arrival, which are crucial for advanced wireless communication systems.

Main advantage of ray-tracing over traditional stochatic models is its ability to provide site-specific realistic channel modeling.
This property enables channel evolves smoothly with user modiblity, making it suitable for evaluating advanced scenarios that are difficult to capture with stochatic models.
For example, ray-tracing can accurately capture the signal drop when a user moves behind a building, which is difficult to model with traditional stochastic approaches.
Furthermore, ray-tracing can provide detailed spatial information about the propagation environment, such as angle of arrival and delay spread, which are crucial for advanced wireless communication systems.

Despite its advantages, ray-tracing is computationally expensive due to the need to evaluate a large number of propagation paths.
Recent advances in computing hardware, especially the use of GPUs and parallel processing techniques, have made ray-tracing more practical for large-scale simulations.